% ============================================
% Section: Cooperative Mapping
% ============================================
\label{sec:mapping}

Each node maintains a local 3DGS map $\GaussianMap_i$ built by the per-node
front-end (Section~\ref{sec:ctfusion}). The cooperative back-end reconciles
these local maps into a globally consistent photo-realistic map through four
mechanisms: submap registration, Gaussian merge, pose-graph optimization, and
loop closure.

\subsection{Local Map and Ownership}
\label{sec:local}

The local map is identical to Gaussian-LIC2's: a 3DGS map densified by a
zero-shot depth completer for LiDAR-blind areas, supervised by LiDAR geometry
and camera photometry, with bounded scale and regularization. The only
addition is that every Gaussian carries an \emph{owner} tag recording the
$\text{node\_id}$ that created it. Ownership enables per-node pruning (a node
leaving the network can offer its Gaussians to the head rather than have them
orphaned) and is the basis for conflict resolution
(Section~\ref{sec:conflict}).

\subsection{Submap Registration}
\label{sec:registration}

When a head receives a submap $\submap_j$ from node $j$ expressed in $W_j$, it
must estimate the aligning transform $T_{W_h \leftarrow W_j}$ into the head's
frame. Three strategies are tried in escalating cost:

\paragraph{(1) Descriptor prior.} If a coarse prior is available (GNSS, prior
map, or initial handshake), it seeds the next step.

\paragraph{(2) Gaussian-to-Gaussian ICP.} Register the incoming submap's
Gaussian means against the head's current merged map using point-to-plane ICP,
where the plane is the local surface normal estimated from the neighboring
Gaussians' scale covariances. Each correspondence is weighted by the product of
the two Gaussians' opacities and the planarity score
$\kappa = \frac{\lambda_3 - \lambda_1}{\lambda_3}$ (eigenvalues of the
Gaussian's scale covariance). Converges within 10 iterations when spatial
overlap exceeds 30\%.

\paragraph{(3) Feature-based fallback.} When overlap is below 30\%, we fall
back to 2D feature matching between representative keyframes of the two nodes
(stored as image descriptors in the submap header), recover $T$ via PnP +
RANSAC, then refine with step (2).

The output is an estimate $T_{W_h \leftarrow W_j}$ with an information matrix
$\Omega$ derived from the ICP Hessian, emitted as an inter-node edge in the
pose graph (Section~\ref{sec:pgo}).

\subsection{Gaussian Map Merge}
\label{sec:merge}

Once aligned, incoming Gaussians are merged into the head's map:

\begin{enumerate}
  \item \textbf{Spatial dedup.} For each incoming Gaussian $g_{\text{in}}$,
    query the map octree for existing Gaussians within
    $\tau_d = 3 \cdot \bar{s}$, where $\bar{s}$ is the mean scale. If a
    candidate $g_{\text{ex}}$ is found, \emph{fuse} the pair: the merged mean
    is the opacity-weighted average, the scale and SH coefficients are
    opacity-weighted averages, and the opacity is the sigmoid of the
    logit-sum (mirroring the consolidation in GS-LIVM
    \cite{xie2025gslivm} and LIV-GS \cite{xiao2025livgs}). If no candidate is
    found, insert $g_{\text{in}}$.
  \item \textbf{Ownership bookkeeping.} Merged Gaussians carry a list of
    contributing $\text{owner}$ ids, so a node can later prune only its own
    contributions.
  \item \textbf{Map delta broadcast.} The head emits a MAP\_DELTA containing
    only the changed Gaussians (new + fused), so member nodes can lazily
    update their local copies for rendering and planning.
\end{enumerate}

\subsection{Pose-Graph Optimization}
\label{sec:pgo}

A robust pose graph is maintained across all nodes:

\begin{itemize}
  \item \textbf{Vertices:} per-node keyframe poses (one per submap close).
  \item \textbf{Edges:}
    \begin{itemize}
      \item intra-node odometry edges, with information matrices from the
        Gauss-Newton Hessian of the CT solve;
      \item inter-node registration edges from Section~\ref{sec:registration};
      \item loop-closure edges from Section~\ref{sec:loop}.
    \end{itemize}
\end{itemize}

The graph is solved by GTSAM's iSAM2 incrementally on the cluster head, and at
coarser granularity on the Tier-0 server. Solutions are pushed down as
MAP\_DELTA corrections; member nodes apply a rigid per-submap correction to
their owned Gaussians (exact for means and rotations; scale and SH are
re-fitted locally over a short window).

\subsection{Loop Closure}
\label{sec:loop}

Two complementary mechanisms:

\paragraph{Geometric.} Submap registration with a low overlap threshold (10\%)
is attempted against graph-distant submaps using the pose graph as a prior;
success creates a loop edge.

\paragraph{Visual.} A NetVLAD descriptor is computed per keyframe and
broadcast to the head. The head maintains a DBoW2-style inverted index and
queries candidate loops when descriptors match within threshold; verified
loops emit PGO constraints.

\subsection{Conflict Resolution}
\label{sec:conflict}

When two nodes observe the same region at different times and produce
divergent Gaussian parameters (e.g., a parked car moved), we resolve the
conflict by a \emph{confidence-weighted} rule:
\begin{equation}
  \text{confidence}(g) = \text{opacity}(g) \cdot \bigl(n_{\text{photo}}(g) + n_{\text{lidar}}(g)\bigr) \cdot w_{\text{sensor}}(\text{owner}(g)),
\end{equation}
where $w_{\text{sensor}}$ is a per-node quality weight from the node
descriptor. Higher-confidence Gaussians dominate the merge; lower-confidence
ones are demoted (opacity shrunk) and eventually pruned. This is a
first-pass treatment of inter-node dynamics; a fuller dynamic-aware extension
along the lines of RoSe-SLAM \cite{wang2026roseslam} and Dynamic-LIVO
\cite{zhang2026dynamiclivo} is left for future work.

\subsection{Algorithm Summary}

\begin{algorithm}[t]
\caption{Per-node cooperative loop (one iteration).}
\label{alg:coop}
\begin{algorithmic}[1]
\State Acquire multi-cam, multi-LiDAR, IMU data with timestamps.
\State Stack IMU, LiDAR, camera, time-offset residuals; solve sliding-window
       Gauss-Newton on the B-spline. \Comment{Section~\ref{sec:ctfusion}}
\State Update local Gaussian map $\GaussianMap_i$; densify blind areas.
\State \textbf{if} submap boundary reached \textbf{then}
  \State \quad Extract spatial sub-volume; compress under budget.
  \State \quad Publish SUBMAP\_PUSH to cluster head.
\State \textbf{end if}
\State Receive incoming submaps and MAP\_DELTA from head.
\State Align incoming submaps (Sec.~\ref{sec:registration}); merge
       Gaussians (Sec.~\ref{sec:merge}); emit PGO constraints.
\State Apply PGO corrections to owned Gaussians.
\State Render \& supervise next frame.
\end{algorithmic}
\end{algorithm}
